\chapter{Background and Related Work}
\label{ch:related-work}

\section{Skill extraction and skill inventories}

Skill extraction identifies capability mentions in unstructured text. It is related to named-entity recognition, but the boundary of a skill can be long, compositional, and context dependent. A phrase such as ``experience with scalable distributed systems'' may include a tool, a concept, an experience qualifier, or several nested skills depending on the annotation rules. SkillSpan addresses this problem by annotating hard and soft skills at span level across 14.5K English job-posting sentences with more than 12.5K annotated spans~\cite{zhang2022skillspan}. It is particularly relevant here because the local export maintains the hard/soft distinction rather than collapsing all skills into one universal class.

Earlier systems connected extraction with normalization. The SKILL system treated identification and normalization as a linked problem at scale~\cite{zhao2015skill}; later work considers large-scale occupational normalization and extreme multi-label retrieval from job descriptions~\cite{hoang2018largescale,bhola2020retrieving}. This distinction is still important. Extraction answers whether a source says ``AWS,'' whereas normalization decides whether a surface form maps to a canonical concept such as Amazon Web Services. Conflating the two can hide an extraction error behind a plausible identifier. SkillLLM\_v2 adopts an explicit separation: its linker is attach-only, so canonicalization should not affect span-level extraction F1.

The Skill-LLM lineage demonstrates that general-purpose LLMs can be adapted for skill extraction~\cite{herandi2024skillllm}. Related job-market research explores LLM extraction, synthetic job-posting generation, and nearest-neighbour occupational skill extraction~\cite{nguyen2024rethinking,magron2024jobskape,zhang2024nnose}. These efforts motivate a generative JSON interface but do not remove the need for controlled local evaluation. In this report, the substantive model claim is limited to the stored Stage-2 Mistral results on the local data path, not a direct benchmark claim against papers using different pre-processing, model versions, or splits.

\section{Named-entity recognition and span modelling}

Sequence tagging is a useful baseline for extraction. BiLSTM-CRF and character-aware variants established effective NER architectures before transformer encoders~\cite{huang2015bilstmcrf,lample2016ner,ma2016bilstmcnncrf}. BERT made contextual token representations broadly available for downstream classification and sequence labelling~\cite{devlin2019bert}; RoBERTa and XLM-R extended the pre-training design space~\cite{liu2019roberta,conneau2020xlmr}. A BERT baseline is consequently an important control in this project, even though its local artifact is explicitly a CPU-fast baseline and does not match the Mistral result.

Span-specific methods emphasise that boundaries, not only token labels, determine correctness. SpanBERT pre-trains representations that focus on contiguous spans~\cite{joshi2020spanbert}, and SpanNER formulates recognition through span prediction~\cite{fu2021spanner}. DyGIE++ extends contextual span representations to entities, relations, and events~\cite{wadden2019dygiepp}. Such work is conceptually aligned with SkillSpan evaluation: a semantically reasonable phrase is still incorrect under exact matching when either boundary or label differs. SkillLLM\_v2 therefore reports exact, partial, token, hard-skill, and soft-skill metrics rather than relying on a single aggregate alone.

GLiNER is especially relevant as an auxiliary source because it supports generalist, open-label entity recognition with a bidirectional transformer~\cite{zaratiana2023gliner}. In the audited repository, however, the standalone GLiNER result is zero-shot/few-shot; fine-tuning was skipped. The thesis treats it as a recall-oriented candidate source and an honest baseline, not as evidence of a trained competitor. Generative approaches such as GPT-NER show that LLMs can produce entity structures directly~\cite{wang2023gptner}. Their principal risk in an extraction pipeline is not merely an invalid JSON string; it is a fluent string that fails to align to the input or rewrites a source span. The alignment and validation code in SkillLLM\_v2 address this operational issue.

\section{Domain adaptation and parameter-efficient fine-tuning}

Domain-adaptive pre-training and fine-tuning are common ways to transfer general language representations to specialised text~\cite{gururangan2020dontstop,howard2018ulmfit}. Full model adaptation is expensive for a large causal model, motivating parameter-efficient methods. LoRA freezes the base weights and introduces trainable low-rank update matrices~\cite{hu2022lora}. QLoRA combines adapter learning with quantized base-model loading and provides the NF4 and memory-management techniques that make 4-bit adaptation practical~\cite{dettmers2023qlora}. These papers motivate the compute-conscious method; they do not supply performance evidence for this particular model or corpus.

SkillLLM\_v2 uses a Mistral Small 3.1 24B Instruct path, a 4-bit configuration, and LoRA settings stored in YAML. The Mistral model family established a widely used open-weight architecture line~\cite{jiang2023mistral}, while the project-specific training summary supplies the more important claim here: a recorded 80-epoch, two-process run with persisted adapter and evaluation artifacts. The thesis differentiates this implemented training path from the broad PEFT literature. It does not assert that all hardware settings or all configurations from the cited work were reproduced.

\section{Taxonomy linking and retrieval}

Taxonomies make a skill representation more useful for aggregation, reporting, and matching. ESCO provides a semantic-interoperability framing for job matching~\cite{levrang2014esco}; ESCOXLM-R and entity-linking work show how job-market language can be related to structured inventories~\cite{zhang2023escoxlmr,zhang2024entitylinking}. Large language models have also been investigated as zero-shot ESCO skill matchers~\cite{clavie2023batteries}. A local implementation cannot be assumed to inherit the coverage or accuracy of those published systems. The audited Stage-2 code provides loaders and aliases for multiple sources, but its reported path is a custom alias-backed linker. Its 0.890 values are linked-candidate coverage, while a supervised canonical-link test set is absent.

Retrieval is relevant at two scopes. Dense passage retrieval~\cite{karpukhin2020dpr}, Sentence-BERT~\cite{reimers2019sentencebert}, SimCSE~\cite{gao2021simcse}, and ColBERT~\cite{khattab2020colbert} represent well-established semantic or late-interaction approaches. Retrieval-augmented generation combines retrieved evidence with a generator in knowledge-intensive tasks~\cite{lewis2020rag}. The project's `ContextualRetriever` is deliberately much smaller: it scores lexical term overlap over contextualised records. It serves a deterministic smoke/demo path and an interface boundary for future dense retrieval. This distinction prevents the report from labelling a lexical overlap class as a complete RAG evaluation.

\section{Graphs, matching, and contextual career intelligence}

Career systems need relationships that simple extraction does not express: a user has a skill; a job requires it; a course teaches it; a missing skill explains a weak match. Skills2Job demonstrates graph-oriented job recommendation~\cite{giabelli2021skills2job}. Person--job fit has also been studied through domain-adapted deep networks and context-aware transformer matching~\cite{bian2019domain,li2020competence}. These studies motivate the graph and matching interfaces in SkillLLM\_v2, but no paper can validate the repository's deterministic weights. The graph in this codebase is an in-memory NetworkX multigraph, and the matching result is an interpretable weighted formula.

Context management raises a different problem: sending an entire profile history with every query increases cost and can obscure fresh evidence. A compact, structured profile can cache stable, attributable facts while a retriever supplies dynamic job text. SkillLLM\_v2 labels this separation CAG/RAG/graph policy: stable profile summary, aliases, and course summaries are cached; job postings and scraped descriptions remain retrieval targets; relational explanations remain in the graph. This is a concrete design and smoke-tested implementation, not an empirical comparison of cache-augmented generation, RAG, and GraphRAG methods.

\section{Additional methodological perspectives}

Soft-skill requirements are prominent but difficult to operationalise consistently in job text~\cite{sayfullina2018softskill,calanca2019responsible}. Implicit-skill extraction also illustrates why a text mention should not be treated as a direct measure of competence~\cite{gugnani2020implicit}. SkillLLM\_v2 consequently limits its extraction claim to explicit span evidence and reserves proficiency, transferability, and recency for caller-supplied profile fields or future evaluation.

Other NER lines make complementary trade-offs. Contextual string embeddings and pooled contextual embeddings show how character-level and document-context signals can affect entity boundaries~\cite{akbik2018flair,akbik2019pooled}; iterated dilated convolutions focused on fast recognition~\cite{strubell2017idcnn}. Broad surveys distinguish representation, decoder, data, and domain choices in deep NER~\cite{li2022nersurvey}. These results motivate reporting the local BERT artifact as a control rather than as an exhaustive representation of token-classification approaches.

Parameter-efficient adaptation also includes prefix tuning and prompt tuning~\cite{li2021prefixtuning,lester2021prompttuning}. Unified text-to-text modelling demonstrates the value of treating NLP outputs as controlled generation~\cite{raffel2020t5}, while BART provides a denoising encoder--decoder perspective~\cite{lewis2020bart}. SkillLLM\_v2 uses neither prefix tuning nor a BART-style architecture; these references clarify why the report calls its method causal-language-model SFT with LoRA rather than treating every structured-output model as equivalent.

The broader open-weight ecosystem includes mixture-of-experts Mistral-family work~\cite{jiang2024mixtral} and reproducibility-oriented fine-tuning tooling such as LlamaFactory~\cite{zheng2024llamafactory}. Synthetic instruction construction is another relevant design direction~\cite{wang2023selfinstruct}. The local Stage-2 experiment uses the SkillSpan-derived JSON targets described in \Cref{ch:methodology}; it does not report a Mixtral comparison, LlamaFactory run, or synthetic-instruction ablation.

Document-link-aware pre-training, such as LinkBERT, provides another perspective on document relationships in language representations~\cite{yasunaga2022linkbert}. In contrast, the current project keeps retrieval, extraction, and graph paths modular: it does not train a link-aware language model or report end-to-end retrieval-grounded generation quality. This modularity is useful for testing because a retrieval error can be distinguished from a source-span error and from a graph-edge error.

The design also treats explanations and correction as part of research quality. Automated annotation studies show that language-model outputs can be useful but must be assessed against task-specific evidence and error patterns~\cite{gilardi2023annotation}. In SkillLLM\_v2, a source span, offset, model source, and canonical attachment are retained so that a future user interface can expose and correct a recommendation's premises. This is a more modest and defensible use of automation than claiming that a generated profile is a verified employment assessment.

\section{Research gap and positioning}

The literature provides strong components: span extraction benchmarks, parameter-efficient adaptation, generalist NER, inventory linking, semantic retrieval, and graph recommendation. The local research gap is their operational composition. A sophisticated wrapper can destroy correct base spans through a boundary edit, a generic verifier rejection, or an aggressive merge. Conversely, a carefully scoped normalizer and profile builder can add useful downstream structure without claiming a span-F1 gain. SkillLLM\_v2 is positioned as an evidence-preserving integration study. Its novel practical thesis is not ``more modules are better''; it is that a primary extractor should be an explicit anchor and that every downstream alteration must be attributable, testable, and rejectable.
